% ════════════════════════════════════════════════════════════════════════════════
\section{Introduction}
\label{sec:intro}

This report documents our group project for the
\emph{Identity and Access Management} (IAA) course at the University of
Aveiro.
The project brief asks us to design
and implement \emph{``a secure user authentication protocol for WireGuard
VPN''}~\cite{donenfeld2017wireguard} together with the supporting client
and server applications, MFA, RBA, an IdP, a custom PAM module, and a
helpdesk enrolment interface.
This report describes what we built, the design decisions behind each
component, and where in the codebase each project requirement is
satisfied.

\subsection{Problem Statement}
WireGuard tunnels packets between peers that already know each other's
Curve25519 public key, private tunnel IP, and pre-shared key.
For a host-to-gateway deployment serving many users, the gateway must
pre-configure all six per-peer parameters before any client can connect,
and clients must avoid private IP collisions out-of-band.
We solve this by building a login layer on top of WireGuard: the user
authenticates through a browser-based MFA flow gated by a risk engine,
and only on success does the gateway generate and hand over the
per-session WireGuard parameters so the tunnel can come up.
The entire stack, login client, gateway daemon, identity provider, PAM
module, and helpdesk, runs from a single \texttt{docker compose up},
with the gateway's WireGuard public key as the only trust anchor.

\subsection{Mapping to Project Requirements}
Table~\ref{tab:compliance} maps every mandatory project requirement to
the section of this report and the file in the codebase that implements
it.

\begin{table}[H]
\centering
\caption{Project requirements vs.\ where they are implemented.}
\label{tab:compliance}
\renewcommand{\arraystretch}{1.15}
\scriptsize
\begin{tabular}{p{3.0cm}p{1.2cm}p{3.2cm}}
\toprule
\textbf{Requirement} & \textbf{Section} & \textbf{Code location} \\
\midrule
UserAuthn application                        & \ref{sec:protocol} & \texttt{client/userauthn.py} \\
SrvAuthn application                         & \ref{sec:protocol} & \texttt{server/srvauthn.py} \\
Six WG parameters per side                   & \ref{sec:wgparams} & \texttt{server/srvauthn.py}, \texttt{common/db.py} \\
WG public key as source of trust             & \ref{sec:trust}    & \texttt{common/noise.py} \\
IdP (OAuth 2.0)                              & \ref{sec:protocol} & \texttt{idp/app.py} \\
HTTP-over-secure-tunnel to IdP               & \ref{sec:proxy}    & \texttt{client/proxy.py}, \texttt{server/proxy.py} \\
OAuth Device Authorization Grant             & \ref{sec:protocol} & \texttt{idp/app.py} \\
MFA: password (challenge-response)           & \ref{sec:mfa-pwd}  & \texttt{idp/app.py}, \texttt{idp/static/js/} \\
MFA: TOTP                                    & \ref{sec:mfa-totp} & \texttt{idp/app.py}, \texttt{common/db.py} \\
MFA: FIDO2                                   & \ref{sec:mfa-fido} & \texttt{idp/app.py} \\
RBA (IP, ASN, failures, time)                & \ref{sec:rba}      & \texttt{idp/rba.py} \\
PAM module orchestrating MFA+RBA             & \ref{sec:pam}      & \texttt{idp/pam\_iaa/src/pam\_vpnmfa.c} \\
Helpdesk enrolment interface                 & \ref{sec:helpdesk} & \texttt{idp/helpdesk/app.py} \\
\bottomrule
\end{tabular}
\end{table}

\subsection{Report Structure}
Section~\ref{sec:architecture} gives the system architecture and
trust model.
Section~\ref{sec:protocol} walks through the authentication protocol
end-to-end and shows how the six WireGuard parameters per side are
assigned.
Section~\ref{sec:mfa} covers the three MFA factors and the helpdesk
enrolment interface.
Section~\ref{sec:pam} describes the custom C PAM module.
Section~\ref{sec:rba} describes the RBA engine.
Section~\ref{sec:security} covers the additional security hardening
measures we implemented on top of the core functionality.
Section~\ref{sec:evaluation} demonstrates that the system works end-to-end.
Section~\ref{sec:conclusion} concludes and notes what is intentionally not
implemented.
